\documentclass[10pt,a4paper]{amsart}
\usepackage[margin=19mm]{geometry}
\usepackage{amsmath,amssymb,mathtools}
\usepackage[scr=rsfs,cal=euler]{mathalfa}
\usepackage{xfrac}
\usepackage[unicode,hidelinks,bookmarks=false]{hyperref}
\newcommand{\ie}{i.e.\ }
\newcommand{\cf}{cf.\ }
\newcommand{\resp}{respectively\ }
\newcommand{\Subsection}{\subsection}
\providecommand{\nobreakdash}{\nobreak\hbox{-}}
\setlength{\emergencystretch}{2em}
\multlinegap0pt
\usepackage{bm}
\usepackage{bbm}
\usepackage{dsfont}
\usepackage{xspace}
\usepackage{cancel}
\usepackage{mathtools}
\usepackage[shortlabels]{enumitem}
\usepackage{amsthm}
\makeatletter
\@ifundefined{theorem}{\newtheorem{theorem}{Theorem}}{}
\@ifundefined{lemma}{\newtheorem{lemma}[theorem]{Lemma}}{}
\makeatother
\newcommand{\makeset}[2]{\left\lbrace #1 \;\middle|\;
 \begin{tabular}{@{}l@{}}
   #2
  \end{tabular}
  \right\rbrace}
\title[A note on intrinsic topologies of groups]{A note on intrinsic topologies of groups}
\author{S. Bardyla, L. Elliott, J. D. Mitchell,, Y. P\'eresse}
\date{Source: arXiv:2506.11500v2 (CC BY 4.0)}
\begin{document}
\maketitle

\noindent\emph{Source and licence.} A note on intrinsic topologies of groups. Version arXiv:2506.11500v2 (CC BY 4.0), distributed under the Creative Commons Attribution 4.0 International licence, as recorded in the arXiv licence metadata for this exact version and shown on the abstract page https://arxiv.org/abs/2506.11500v2. Full rights notice: licenses/arxiv-2506.11500-CC-BY-4.0.txt.\par\smallskip

\noindent\textit{Editorial scope.} Counted unit: the Lemma whose source label is \texttt{SN nice zariski neighbourhood} in the article (source0.tex lines 427--435, source label \texttt{SN nice zariski neighbourhood}), delivered together with the article's own complete original proof of it (source0.tex lines 436--487, one \texttt{proof} environment of 52 lines). No mathematics is written, repaired or re-derived: statement and proof are the article's own text line for line. Required context retained uncounted: none. Every definition, display and label of those lines is retained unchanged. Omitted from this package and not redistributed: 1--426, 488--1094. Nothing inside a retained excerpt is omitted or altered: every retained body is delivered exactly as the article's lines are written, so no omission receipt of any kind accompanies this package. Citations: the retained excerpts carry no citation command at all, so no bibliography entry of the article is transcribed and no reference is redistributed. Reference adaptation: none was needed and none was made; every reference inside the delivered unit resolves inside it, so no unresolved reference or citation is produced. Printed slips kept literal: no printed word, symbol, hypothesis or number of the counted statement or of its proof was altered. Layout and class: the article's own class is \texttt{amsart} with packages including color, amssymb, thm-restate, hyperref, ulem, enumitem, cleveref, xy, graphicx; the wrapper is the lane's pinned \texttt{amsart} editorial wrapper at 10pt on A4, which restates the theorem-like environments the retained text needs and adds the article's own macro definitions that the retained text needs (\texttt{\textbackslash makeset}), with extra wrapper packages (bm, bbm, dsfont, xspace, cancel, mathtools, enumitem). The delivered numbering is the unit's own. Assets: no retained span contains a figure environment, a graphics-inclusion command, a PDF-page inclusion, a table or a TikZ picture; no image file is copied into this package, the package declares no asset dependency and no image file is redistributed with it. Licence: version arXiv:2506.11500v2 of this article is distributed under the Creative Commons Attribution 4.0 International licence, as recorded in the arXiv licence metadata for this exact version and shown on the abstract page https://arxiv.org/abs/2506.11500v2. A full rights notice is delivered at licenses/arxiv-2506.11500-CC-BY-4.0.txt. Characters: every retained excerpt is pure ASCII, so the delivered engine, XeLaTeX with its default Latin Modern text font, reports no missing character.
\begin{lemma}\label{SN nice zariski neighbourhood}
Let \(S\) be a cancellative monoid.
The semigroup Zariski topology on \(S\) has a basis consisting of the sets \(N_{A, B}\) where \(A, B\) are ragged matrices over \(S\) such that
\begin{enumerate}
    \item \(d_{A,i} > 0\) or \( d_{B,i}>0\),
      \item \(a_{i, 0} \neq b_{i, 0}\),
\end{enumerate}
for every $i\in k$ where \(k\) is the number of rows of \(A\) and \(B\).
\end{lemma}

\begin{proof}
By the definition of the semigroup Zariski topology, the nonempty sets of the form \(N_{A,B}\), where \(A,B\) are ragged matrices with the same amount of rows, form a basis in $(S,\mathfrak{Z}(S))$. We show that if \(N_{A,B}\neq \varnothing\), then \(N_{A, B}=N_{P, Q}\) for some \(P, Q\) satisfying conditions (1) and (2).

For each pair of ragged matrices \((A, B)\) over \(S\) with \(k\) rows each, let \(t_{A, B}\) be the tuple \[(k, d_{A,0}, \ldots, d_{A,k-1}, d_{B,0}, \ldots, d_{B,k-1}).\] 
For a given pair of $k$-rowed ragged matrices $(A,B)$ such that $N_{A,B}\neq\varnothing$, let $\mathcal{PQ}$ be the set of all pairs \((P, Q)\) of $k$-rowed ragged matrices such that \(t_{P, Q}\) is the minimum element of
\(\makeset{t_{P, Q}}{ \(N_{P, Q}= N_{A, B}\)}\),
where we order these tuples lexicographically.

% We may assume that \(k\in \omega\) is minimal with this property and for all \(i\in k\) there are no \(A'\) and \(B'\) each with \(k\) rows, \(N_{A, B}=N_{A', B'}\), \(A\) and \(A'\) agree in all rows other than \(i\), \(d_i(A')<d_{i}(A)\).
% \\
% Similarly with \(B\).

We claim that each pair $(P,Q)\in\mathcal{PQ}$ satisfies condition (1). 
Suppose to the contrary that \(d_{P,i}=d_{Q,i}=0\) for some \(i\in k\).
As \(N_{P, Q}\) is not empty, we have \(p_{i, 0}\neq q_{i, 0}\).
Thus
\(N_{P, Q} =N_{A, B}\)
where \(A, B\) are obtained from \(P, Q\) by deleting the \(i^{th}\) row, 
which contradicts the minimality of \(k\).


Let us check that there exists a pair $(P',Q')\in\mathcal{PQ}$ that satisfies condition (2). 
Fix any $(P,Q)\in \mathcal{PQ}$. If \(p_{i, 0}\neq q_{i, 0}\) for each $i\in k$, we are done. Otherwise, let $I=\{i\in k: p_{i, 0}= q_{i, 0}\}$. For each $i\in I$ the following three cases require consideration:
\begin{enumerate}[\rm(i)]
\item \(d_{P,i} \neq 0\) and \(d_{Q,i}\neq 0\);
\item \(d_{P,i} = 0\) and \(d_{Q,i}\neq 0\);
\item \(d_{P,i} \neq 0\) and \(d_{Q,i}= 0\).
\end{enumerate}


%\underline{Claim 2}
%If for some \(i\in k\), \(d_{A,i} \neq 0\) and \(d_{B,i}\neq 0\), then \(A[{i, 0}]\neq B[{i, 0}]\).\\
%\underline{Proof of Claim:}\\
%Suppose that \(A[{i, 0}]= B[{i, 0}]\). 
(i) As \(S\) is cancellative for all \(x\in S\),
\[p_{i,0}xp_{i, 1}x \ldots xp_{i, d_{P,i}} \neq q_{i,0}xq_{i, 1}x \ldots xq_{i, d_{Q,i}}\]
if and only if 
\[p_{i, 1}x \ldots xp_{i, d_{P,i}} \neq q_{i, 1}x \ldots xq_{i, d_{Q,i}}.\]
The latter polynomials have smaller degree contradicting the minimality of \(d_{P,i}\) and \(d_{Q,i}\). Hence case~(i) is impossible.

(ii) Let \(f\in S\setminus \{1_{S}\}\).
Then
\[p_{i, 0} \neq Q_{i}(x)\qquad \iff \qquad p_{i, 0}f \neq Q_{i}(x)f.\]

In this case we redefine \(p_{i, 0}:= 
p_{i, 0}f\) and \(q_{i, d_{Q,i}}:= q_{i, d_{Q,i}}f\). 
Note that \(p_{i,0}=p_{i, 0}1_S\neq p_{i, 0}f\) by cancellativity.

Case (iii) is similar to case (ii).

It is clear that the described above process of redefining entries of $P$ and $Q$ gives us a pair $(P',Q')\in\mathcal{PQ}$ that satisfies condition (2).
\end{proof}

\end{document}
